\subsubsection{Subproceso 3.B: Confirmación y Registro Oficial}

La revisión documental y la observación con apuntes de campo permitieron obtener los siguientes resultados sobre el funcionamiento real del Subproceso 3.B:

\textbf{Actividades observadas y flujo efectivo:}
\begin{itemize}
    \item Al recibir la Pre-Ficha Aprobada o Subsanada, el sistema cierra la edición para el estudiante, guarda sus cursos en el SIGA y emite el Reporte de Matrícula de forma automática (TUPA PA5400954D, RR-3698-2024).
    \item El cierre se realiza cuando la Facultad avisa a la DIRCE que terminó su matrícula y le envía el archivo digital con el padrón final (RR-0570-2022, Art. 78\textdegree{}); luego, la Oficina de Estadística de la Facultad firma el Reporte y la Boleta.
    \item Las listas de clase se entregan a los profesores el primer día hábil de la segunda semana (RR-0570-2022, Art. 78\textdegree{}), y el registro de matrícula junto con las boletas se entregan a las Facultades en un máximo de diez días hábiles tras el cierre (RR-0570-2022, Art. 80\textdegree{}).
\end{itemize}

\textbf{Tiempos de espera identificados:}
\begin{itemize}
    \item \textbf{Listas desactualizadas en la primera semana:} Hay alumnos que asisten a clases sin aparecer en las listas provisionales de los profesores hasta que se terminan de registrar las rectificaciones aprobadas (RR-0570-2022, Art. 78\textdegree{}).
    \item \textbf{Demora en la firma de boletas:} El exceso de carga laboral en las Oficinas de Estadística retrasa la entrega de las boletas firmadas que los alumnos necesitan para trámites de becas y seguros.
    \item \textbf{Espera de hasta diez días hábiles tras el cierre:} Entre el cierre de matrícula y la entrega final del registro oficial (RR-0570-2022, Art. 80\textdegree{}), la Facultad y los estudiantes trabajan solo con documentos y listas provisionales.
\end{itemize}

\textbf{Problemas identificados:}
\begin{itemize}
    \item \textbf{La Facultad no tiene acceso en el sistema para modificar matrículas:} Es el problema más grave reportado en esta etapa. Los ajustes de horario acordados con alumnos y docentes no pueden hacerse directamente en la Facultad; todo cambio debe pedirse a la DIRCE, lo que genera trámites manuales y retrasa la actualización de las listas de clase.
    \item \textbf{Aviso de cierre sin revisión conjunta entre Facultad y DIRCE:} La Facultad solo envía un aviso informativo al terminar su matrícula, sin una revisión cruzada con la DIRCE para confirmar que no falten alumnos o cambios por procesar.
    \item \textbf{Retrasos acumulados de las etapas anteriores:} Los problemas no resueltos en los subprocesos previos (deudas no filtradas, fichas creadas con datos del ciclo anterior, horarios cargados tarde y cambios de plan sin convalidar) llegan a la primera semana de clases como matrículas rezagadas y rectificaciones de última hora, obligando a rehacer los padrones.
    \item \textbf{Errores al enviar datos a la SUNEDU:} Las diferencias en algunos datos del padrón retrasan la validación en la SUNEDU y demoran la entrega del carné universitario.
\end{itemize}

En resumen, los problemas de esta etapa final son el resultado acumulado de las fallas anteriores: los trámites pendientes y las rectificaciones tardías se juntan en la semana de cierre y, como la Facultad no puede editar el sistema directamente, cada corrección debe pasar por la DIRCE, alargando la entrega de listas oficiales y de boletas firmadas.
